% 使用 XeLaTeX 编译
\documentclass[12pt,a4paper]{ctexart}
\usepackage{amsmath,amssymb,amsthm}
\usepackage{enumitem}
\usepackage{geometry}
\geometry{left=2.5cm,right=2.5cm,top=2cm,bottom=2cm}

\title{\textbf{群论期中考试试题（回忆版）}}
\author{阮东}
\date{\today}

\begin{document}
\maketitle

\section*{一、概念叙述（33分）}
\begin{enumerate}[label={\arabic*.}]
    \item Lagrange定理
    \item Maschke定理
    \item 酉定理
    \item Burnside定理
    \item 不变子空间
    \item Cayley定理
    \item 商群
    \item 同态核定理
    \item 不变子群
\end{enumerate}

\section*{二、判断简答题（35分）}
\begin{enumerate}[label={\arabic*.}]
    \item 群G的正规子群的正规子群是G的正规子群
    \item 若两个群不同构，则特征标表不同
    \item 一个群能写成子群的半直积的条件是什么，叙述这两个子群的作用
    \item 求表示空间函数基的先决条件是什么
\end{enumerate}

\section*{三、计算题（共33分）}
\begin{enumerate}[label={\arabic*.}]
    \item 证明$f_{a,c}(x) = \frac{ax}{cx+1}$在函数复合运算下成群，计算其子群、正规子群和共轭类。

    \item $D_3$群，取函数基为
    $$\psi_1 = e^{-ix},\ \psi_2 = e^{-i(-\frac{1}{2}x-\frac{\sqrt{3}}{2}y)},\ \psi_3 = e^{-i(-\frac{1}{2}x+\frac{\sqrt{3}}{2}y)},\ $$
    计算群的表示、特征标表，并约化之。
\end{enumerate}
\vfill
\noindent\textbf{注：}本试卷为回忆版，题目及分值可能与原卷存在出入，仅供参考。

\end{document}
